\section{Correlation of hydrate rim measurements with \acs*{EBSD} orientations}
\label{sec:ebsd-hydrate}

%The  analysis  yields, for each valid angular step around a clinker grain, a rim length in the plane of the polished section.
To relate the \ac{HLT} measurements which were measured as described in \zcref{subsec:algo_hlm} to the crystallographic orientation of the same clinker volume, the \ac{BSE} micrographs were combined with \ac{EBSD} data acquired on the identical surface.
%The following sections summarise how the two datasets were registered, how crystal directions were assigned to each rim measurement, and which limitations arise from the two-dimensional character of both modalities.

The carbon coated specimen is therefore mounted on a pre-tilted specimen holder and tilted to a \qty{70}{\degree} angle to the electron beam.
The specimen is then scanned at \qty{15}{\kilo\volt}, \qty{1.6}{\milli\ampere} at a step-size of \qty{2}{\micro\meter\per\px} to acquire the \textsc{Kikuchi} patterns of \ac{C3S} and \ac{CSH}.

The \ac{EBSD} experiment produces a 2D dataset of crystal orientation estimates, phase indices and pattern quality indicators. 
The raw data is preprocessed in \textit{AZtecCrystal 3.2} \cite{azteccrystal.2010} by \textit{Oxford Instruments} to clean up pseudosymmetric solutions.
Finally, the data is exported as channel text file (.ctf) which can be further processed using \textit{Python}.

%Each scan pixel is associated with \textsc{Euler} angles in the \textsc{Bunge} convention, which define a rotation $g$ from the sample reference frame to the crystal reference frame of the indexed phase.

The \ac{BSE} image was manually aligned to the \ac{EBSD} map (using the band-contrast image) in \textit{QGIS 3.42} \cite{QGIS_software} by identifying common points.
The \textit{QGIS}-project was then used to projective transform the high-resolution \ac{BSE} image to the low resolution \ac{EBSD} dataset.

%After alignment, the physical step sizes of the \ac{EBSD} grid and the transformed \ac{BSE} image were harmonised by an integer upscale of the segmentation used in the rim algorithm and a scalar alignment factor between low-resolution and high-resolution \ac{BSE} exports, so that one \ac{EBSD} column and row index corresponds to a consistent footprint on the section.

For every rim measurement, the clinker-hydrate boundary point reconstructed in \ac{BSE} coordinates was projected onto the nearest \ac{EBSD} grid node. % (including segmentation upscale and \ac{ROI} origin)  $(i_x,\,i_y)$.
The \textsc{Euler} angles at that node then represent the orientation attributed to the measurement.
Where several angular steps around one grain map to the same grid node, multiple records can share identical indices.%; this is expected when the \ac{EBSD} step width is large compared with the arc length sampled along the rim.

\color{red}
% TODO: Text umschreiben - hier nur Platzhalter!
In the polar representation of the rim algorithm (\zcref{fig:hlt-method}), each valid direction is described by an azimuth $\theta_\mathrm{idx}$ in the plane of the micrograph (sample surface).
A unit vector $\mathbf{n}_\mathrm{sample}$ was constructed in the sample reference frame with components
\begin{equation}
	\label{eq:nsample-plane}
	\mathbf{n}_\mathrm{sample} = \bigl(\cos\theta,\; \sin\theta,\; 0\bigr)^\mathrm{T},
\end{equation}
where $\theta$ follows from the discrete index and an optional sign convention for the image $y$-axis.

The crystal orientation at the associated grid point yields $\mathbf{g}(\varphi_1,\,\Phi,\,\varphi_2)$ in the \textsc{Bunge} convention ($\mathbf{g} = \mathbf{R}_{z}(\varphi_2)\,\mathbf{R}_{x}(\Phi)\,\mathbf{R}_{z}(\varphi_1)$).
The corresponding rim direction in the crystal reference frame is
\begin{equation}
	\label{eq:h-crystal}
	\mathbf{h} = \frac{\mathbf{g}\,\mathbf{n}_\mathrm{sample}}{\lVert \mathbf{g}\,\mathbf{n}_\mathrm{sample} \rVert},
\end{equation}
with components $(h_x,\,h_y,\,h_z)^\mathrm{T}$.
Because $\mathbf{h}$ is normalised, all measurements lie on the unit sphere in crystal coordinates.

For the cubic clinker phase investigated here, the angle between the measured rim normal direction and the $\langle001\rangle$ family was summarised by
\begin{equation}
	\label{eq:angle-001}
	\alpha_{001} = \arccos\!\bigl(\lvert h_z\rvert\bigr),
\end{equation}
i.e.\ the acute angle between $\mathbf{h}$ and the crystal $c$-axis, using the absolute value of $h_z$ so that opposite poles of the same axis are treated equivalently in the statistical plots.
The scalar rim length from the hydrate analysis (denoted $l$ in the workflow output) was retained alongside $\alpha_{001}$ and $(h_x,\,h_y,\,h_z)$ for joint distributions, spatial maps on $(i_x,\,i_y)$, and spherical heat maps of measurement density on the orientation sphere.

\FloatBarrier
\section{Phase consistency and optional orientation infill}

The phase map delivered with the \ac{EBSD} data does not always label every pixel that appears as residual clinker in the \ac{BSE} segmentation---for example when pattern indexing failed at the grain boundary or when slight misalignment shifts the boundary between pixels.
To reduce missing orientations at valid rim locations, an optional correction was applied on the \ac{EBSD} grid: pixels that fell inside a clinker mask derived from the upscaled \ac{BSE} segmentation but were not indexed as clinker in the \ac{CTF} were assigned the \textsc{Euler} angles of the nearest indexed clinker pixel within the same connected clinker region, subject to a maximum search radius in grid units.
Pixels outside that radius were left unchanged so as not to interpolate across unrelated grains.
Whenever this infill was used, the records were flagged in the exported table to allow filtering by data source (\ac{CTF} original versus infilled).

\FloatBarrier
\section{Multi-site compilation}

Measurements from several fields of view or specimen positions were exported with an explicit site identifier, concatenated into a single table, and analysed with the same graphical procedures as the single-site workflow.
Histograms, scatter plots, and heat maps then represent the pooled angular statistics, whilst maps of $(i_x,\,i_y)$ and interactive orientation plots are evaluated separately per site, because pixel indices are only meaningful within each \ac{EBSD} map.

\FloatBarrier
\section{Limitations}

The procedure links a 2D rim length in the section plane to a 3D crystal orientation.
No information is recovered along the surface normal; the rim vector $\mathbf{n}_\mathrm{sample}$ is confined to the plane of polish.
Furthermore, $\mathbf{g}$ refers to the orientation of the voxel averaged by the \ac{EBSD} step at $(i_x,\,i_y)$; sub-grain orientation spread, pseudosymmetry, and indexing ambiguity are not resolved beyond the vendor indexing output.
Finally, correlation plots between $l$ and $\alpha_{001}$ are descriptive: they do not, by themselves, establish a causal relationship between crystallography and local growth rate without additional modelling of transport and reaction at the three-phase boundary.

\color{black}